% !PRINTED 589
\weberpage{589}
%<<CONT
ist. \gesperrt{Wir nennen $\Psi$ die zum Ideal $\mathfrak{a}$ gehörige Gruppe, und wir sagen auch,
$\mathfrak{a}$ gehört zu der Gruppe $\Psi$.}

Da man in jeder Gleichung zwischen Zahlen und Functionalen in $\Omega$ alle
Substitutionen der Gruppe $\Phi$ ausführen darf, wobei die Grössen des Körpers $R$
ungeändert bleiben, ohne dass die Gleichung zu bestehen aufhört, so folgt, dass
Einheiten, ganze Functionale, associirte Functionale, durch einander theilbare
Functionale diese Eigenschaften nicht verlieren, wenn irgend eine der Substitutionen von
$\Phi$ ausgeführt wird. Wir heben den Satz hervor:

1. \gesperrt{Ist $\mathfrak{p}$ ein Primideal, so sind auch alle mit $\mathfrak{p}$ in
Bezug auf $R$ conjugirten Ideale $\mathfrak{p} \mid \varphi$ Primideale. Ist $\mathfrak{a}$
durch irgend eine Potenz von $\mathfrak{p}$ theilbar so ist $\mathfrak{a} \mid \varphi$
durch die gleiche Potenz von $\mathfrak{p} \mid \varphi$, theilbar.} % [sic] source omits the comma after "theilbar" and prints a stray comma after $\mathfrak p\mid\varphi$

Ist $\mathfrak{p}$ ein Primideal in $\Omega$, so giebt es eine und nur eine durch
$\mathfrak{p}$ theilbare natürliche Primzahl $p$.

Zerlegt man diese in ihre Primfactoren in $R$, so muss einer dieser Primfactoren, den
wir mit $\mathfrak{P}$ bezeichnen, durch $\mathfrak{p}$ theilbar sein.

Ist dann $\mathfrak{A}$ irgend ein Ideal in $R$ und zugleich durch $\mathfrak{p}$
theilbar, so ist der grösste gemeinschaftliche Theiler von $\mathfrak{P}$ und
$\mathfrak{A}$, der nach \S.\,139 auch in $R$ enthalten ist, durch $\mathfrak{p}$
theilbar, und ist also gewiss keine Einheit. $\mathfrak{P}$ und $\mathfrak{A}$ sind also
nicht relativ prim, und $\mathfrak{A}$ ist folglich durch das Primideal $\mathfrak{P}$
theilbar. Ist $\mathfrak{A}$ auch ein Primideal, so müssen $\mathfrak{A}$ und
$\mathfrak{P}$ identisch sein. Damit ist, mit Rücksicht auf 1., bewiesen:

2. \gesperrt{Ist $\mathfrak{p}$ ein Primideal in $\Omega$, so giebt es ein und nur ein
Primideal $\mathfrak{P}$ in $R$, das durch $\mathfrak{p}$ theilbar ist, und jedes durch
$\mathfrak{p}$ theilbare Ideal in $R$ ist zugleich durch $\mathfrak{P}$ theilbar. Das
Primideal $\mathfrak{P}$ ist zugleich durch die sämmtlichen zu $\mathfrak{p}$
conjugirten Primideale $\mathfrak{p} \mid \varphi$ theilbar.}

Die Partialnorm $\mathfrak{N}_R(\mathfrak{p})$ ist durch $\mathfrak{p}$, und folglich als
Ideal in $R$ durch $\mathfrak{P}$ theilbar. Sie kann aber auch nach dem Satze 1. durch
kein von $\mathfrak{P}$ verschiedenes Ideal in $R$ theilbar sein. Denn ist $\mathfrak{P}'$
irgend ein Primtheiler von $\mathfrak{N}_R(\mathfrak{p})$, so ist $\mathfrak{P}'$
wenigstens durch einen der Primfactoren $\mathfrak{p} \mid \varphi$ theilbar und mithin
gleich $\mathfrak{P}$. Folglich
%CONT>>
